\documentclass[10pt,a4paper]{amsart}
\usepackage[margin=19mm]{geometry}
\usepackage{amsmath,amssymb,mathtools}
\usepackage[scr=rsfs,cal=euler]{mathalfa}
\usepackage{xfrac}
\usepackage[unicode,hidelinks,bookmarks=false]{hyperref}
\newcommand{\ie}{i.e.\ }
\newcommand{\cf}{cf.\ }
\newcommand{\resp}{respectively\ }
\newcommand{\Subsection}{\subsection}
\providecommand{\nobreakdash}{\nobreak\hbox{-}}
\setlength{\emergencystretch}{2em}
\multlinegap0pt
\usepackage{mathtools}
\usepackage[shortlabels]{enumitem}
\usepackage{amssymb}
\usepackage{tikz}
\newtheorem{prop}{Proposition}
\newcommand{\C}{\mathfrak{C}^{\,}}
\newcommand{\inv}{\operatorname{inv}}
\title{On generalized morphisms associated with \\ endofunctors of $C^*$-algebras}
\author{Georgii S. Makeev}
\date{Source: arXiv:2509.02001v2 (CC BY 4.0)}
\begin{document}
\maketitle

\noindent\emph{Source and licence.} On generalized morphisms associated with \\ endofunctors of $C^*$-algebras. Version arXiv:2509.02001v2 (CC BY 4.0), distributed by arXiv under the Creative Commons Attribution 4.0 International licence, http://creativecommons.org/licenses/by/4.0/, as recorded in the arXiv licence metadata for this exact version and shown on the abstract page https://arxiv.org/abs/2509.02001v2. Full licence notice: \texttt{licenses/arxiv-2509.02001-CC-BY-4.0.txt}.\par\smallskip

\noindent\textit{Editorial scope.} Counted unit: the article's prop prop (source0.tex lines 4107--4113). The article's complete original proof of it is delivered with it (source0.tex lines 4115--4172, one \texttt{proof} environment of 58 lines). No mathematics is written, repaired or re-derived: the statement and the proof are the article's own lines, in the article's own notation, and no retained line is substituted or re-flowed.  No further source-local context is needed: every cross-reference of every retained line resolves inside the two delivered spans.  Nothing was omitted from any retained span, so the package carries no omission record.  No retained line contains a citation, so the wrapper carries no bibliography and no bibliography entry.  No retained line contains a figure environment, an image-inclusion command or drawing code, so no image is delivered and the package declares no asset dependency.  Editorial formatting only, in addition: an AMS \texttt{amsart} preamble that restates the article's own declarations and macros used by the retained lines, namely the environments \texttt{prop} on separate counters, together with the standard packages those lines need.  The retained environments print in this document as Proposition 1 in order of appearance, whereas the article prints the same items with the numbering of its own full document.  The complete native source of the article as distributed in the arXiv e-print for this exact version is bundled unchanged as \texttt{sources/184289/source0.tex}.
\begin{prop}
The suspension functor $S\coloneqq\C_{(0,1)}$ is homotopy symmetric.
Specifically, one can define $\inv_{S}$ by the formula
\[
\inv_{S}B\colon C_{0}((0,1),B)\to C_{0}((0,1),B)\colon f\mapsto(t\mapsto f(1-t)).
\]
\end{prop}

\begin{proof}
For $0\leq a\leq b\leq1$ and $f\in C_{0}((0,1),B)$,
let us define
\[
T_{a,b}(f)(t)\coloneqq\begin{cases}
f\left(\frac{t-a}{b-a}\right), & t\in[a,b]\\
0, & \textrm{otherwise}.
\end{cases}
\]
For all $f,g\in C([0,1],B)$, we define the functions $\overline{f},f*g\in C([0,1],B)$
by the formulas
\begin{alignat*}{3}
 & \overline{f}(t)\coloneqq f(1-t), & \qquad\qquad\qquad\qquad & f*g(t) & \coloneqq\begin{cases}
f(2t), & t\in\left[0,\frac{1}{2}\right];\\
g(2t-1), & t\in\left[\frac{1}{2},1\right].
\end{cases}
\end{alignat*}
The statement of the proposition follows from the chain of homotopies
\[
\left[f\mapsto\left(\begin{array}{cc}
f & 0\\
0 & \overline{f}
\end{array}\right)\right]\simeq\left[f\mapsto\left(\begin{array}{cc}
T_{0,\frac{1}{3}}(f) & 0\\
0 & T_{\frac{2}{3},1}(\overline{f})
\end{array}\right)\right]\simeq\left[f\mapsto\left(\begin{array}{cc}
T_{0,\frac{1}{3}}(f)*T_{\frac{2}{3},1}(\overline{f}) & 0\\
0 & 0
\end{array}\right)\right]\simeq0
\]
given by the following formulas:
\begin{alignat*}{2}
 & H_{1}(f)(s)\coloneqq &\ 
& \left(\begin{array}{cc}
T_{0,\frac{1}{3}s+(1-s)}(f) & 0\\
0 & T_{\frac{2}{3}s,1}(\overline{f})
\end{array}\right),\\
 & H_{2}(f)(s)\coloneqq &\ 
& V_{s}\left(\begin{array}{cc}
T_{0,\frac{1}{3}}(f) & 0\\
0 & T_{\frac{2}{3},1}(\overline{f})
\end{array}\right)V^{*}_{s},\\
 & H_{3}(f)(s)\coloneqq &\ 
& \left(\begin{array}{cc}
T_{0,\frac{1}{3}+\frac{s}{1-s}}(f)*T_{\frac{2}{3}-\frac{s}{1-s},1}(\overline{f}) & 0\\
0 & 0
\end{array}\right),
\end{alignat*}
where $s\in[0,1]$ is the homotopy variable, and where $s\mapsto V_{s}$
is a continuous path of unitaries such that $V_{s}=\left(\begin{array}{cc}
1 & 0\\
0 & 1
\end{array}\right)$ for $s\in\left[0,\frac{1}{3}\right]$ and $V_{s}=\left(\begin{array}{cc}
0 & 1\\
-1 & 0
\end{array}\right)$ for $s\in\left[\frac{2}{3},1\right]$.

\end{proof}

\end{document}
